\section{서론}\label{sec:intro}

스파이킹 신경망(Spiking Neural Network, SNN)은 뉴런 간에 이진
스파이크로 정보를 교환하는 신경망 모델로\cite{maass1997networks},
이벤트 기반 연산의 저전력 잠재력 때문에 엣지 추론의 후보로 주목받아
왔다\cite{roy2019towards}. 저자들의 선행 연구\cite{cha2026fusion}는
전용 뉴로모픽 하드웨어\cite{davies2018loihi,merolla2014truenorth}가
아닌 범용 MCU(STM32F746, Cortex-M7 200MHz) 위에서, 표준 TinyML
스택---TensorFlow Lite for Microcontrollers(TFLM)\cite{david2021tflm}와
CMSIS-NN\cite{lai2018cmsisnn}---에 Conv-AvgPool-LIF 퓨전 커스텀
연산자를 통합하여 CIFAR-10 이미지당($T{=}30$ 타임스텝) $2{,}407.0$ms의
추론 지연시간을 실측하였다. 비퓨전 표준 경로 대비 4.0배의 가속이었다.

그러나 이 수치도 하드웨어의 상한과는 거리가 멀다. 퓨전 커널의 유효
처리량은 약 70MMAC/s로, 이 칩의 SMLAD 이론 피크(사이클당 2 MAC,
400MMAC/s)의 약 18\%에 그친다. 본 논문은 이 5.5배의 간극이 어디서
오는지를 보드 실측으로 재진단하고, 그 결과를 근거로 커널을 다시
설계하여 간극의 대부분을 회수한 기록이다.

재진단의 결론은 두 가지였고, 둘 다 통념과 달랐다. 첫째, \textbf{명령
페치는 병목이 아니었다.} 사전 하드웨어 조사는 이 보드가 ART
가속기가 꺼진 채 ITCM 별칭 주소에서 부팅함을 근거로 페치 경로의 큰
손실을 예측했으나, 절제(ablation) 실측에서 ART 활성화와 핫 커널의
ITCM-RAM 배치는 각각 측정 노이즈 이내($\le$0.12\%)의 변화만 보여
개선 효과가 없었다. 코드는 AXIM 경로에서
실행되어 I-캐시가 페치를 담당하고 있었고, 핫 루프는 4KB I-캐시에
충분히 들어가는 크기였다(III장). 둘째, 병목은
데이터 로드였는데, 그 로드의 대부분은 \textbf{0과의 곱셈을 위한
것}이었다. 저자들의 호스트 실측(테스트셋 10{,}000장 기준,
\ref{sec:sparsity}절)에서 이 모델의 발화율은 입력 47.7\%, 첫 합성곱
출력(LIF1) 4.1\%로, 전체 MAC의 84\%를 차지하는 두 번째 합성곱에서는
연산의 약 96\%가 값이 0인 입력에 대한 곱셈-누산이었다.

SNN의 활성값이 이진 $\{0, S\}$라는 사실은 내적을 대수적으로
붕괴시킨다: $\sum_k w_k x_k = S\cdot\sum_{x_k\neq 0} w_k$. 즉 발화
위치의 가중치를 정수로 합산하고 스파이크 진폭 $S$를 마지막에 한 번
곱하면, \textbf{곱셈 없이} dense 경로와 비트 단위로 동일한 결과를
얻는다. 본 논문은 이 항등식을 Cortex-M7의 메모리 계층에 맞게
구현한 희소 스파이크 구동(spike-driven) 커널을 제안한다. im2col을
실체화하지 않고 수용 영역에서 비영 인덱스 목록을 직접 만들고, 전치
가중치 레이아웃 $[k][C_o]$와 8채널 레지스터 블로킹으로 가중치
gather의 캐시 지역성과 레지스터 상주를 동시에 확보하며, requantize와
LIF 에필로그는 dense 경로와 같은 함수$\cdot$같은 순서를 유지한다.
제안 커널은 TFLM 커스텀 연산자 프레임 안에 그대로 통합되어 CMSIS-NN
지원 함수(\texttt{arm\_nn\_requantize\_s64} 등)를 재사용한다.

본 논문의 기여는 다음과 같다.
\begin{itemize}
\item \textbf{실측 기반 병목 재진단}: ART$\cdot$ITCM$\cdot$DTCM$\cdot$커널을
  누적 절제한 5개 빌드 변형의 사이클 정밀 실측으로, 페치 최적화
  2종이 무효과임을 보이고 병목을 데이터 로드로 특정한다. 부팅 주소와
  실행 주소의 불일치라는, 사전 조사가 놓친 원인을 규명한다.
\item \textbf{무곱셈 희소 커널}: 이진 활성의 합 항등식을 이용해
  MAC을 제거하고, im2col 제거$\cdot$가중치 전치$\cdot$8채널 레지스터
  블로킹으로 Cortex-M7에서 동일 메모리 배치의 dense 대비 두 번째
  합성곱 13.0배를 얻는, 구성상(by construction) 비트 동일한 커널을
  제시한다.
\item \textbf{5회 반복의 실리콘 검증 기록}: 커널을 한 번에 설계하지
  않고 보드 실측으로 5회 반복 개선했으며, 그중 한 번의 실측 퇴행(가중치
  전치가 DTCM 상주 레이어에서 역효과)과 그 원인$\cdot$해결까지 보고한다.
\item \textbf{종단 결과와 검증}: 기준 재현 $2{,}328.5$ms에서
  $380.6$ms로 \textbf{6.12배} 가속하면서, 두 실행 경로(TFLM$\cdot$직접
  호출)를 포함한 여섯 빌드 변형의 출력 스파이크 체크섬이 전부 비트
  일치함을 보인다. TFLM 런타임 오버헤드는 2.1\%로 유지된다.
\end{itemize}

이후 구성은 다음과 같다. II장은 대상 모델과 기존
구현을 요약한다. III장은 절제 실측으로 병목을
재진단한다. IV장은 희소 커널의 설계와 반복 개선 과정을,
V장은 종단 실험과 검증을 제시한다. VI장에서 결론을 맺는다.
